\documentclass[11pt,a4paper]{article}
\usepackage[utf8]{inputenc}
\usepackage[T1]{fontenc}
\usepackage[brazil]{babel}
\usepackage{geometry}
\usepackage{graphicx}
\usepackage{booktabs}
\usepackage{longtable}
\usepackage{array}
\usepackage{float}
\usepackage{xcolor}
\usepackage{listings}
\usepackage{pdflscape}
\usepackage{placeins}
\usepackage[hidelinks]{hyperref}
\geometry{margin=2.1cm}
\setlength{\parskip}{0.5em}
\setlength{\parindent}{0pt}
\lstset{
  basicstyle=\ttfamily\footnotesize,
  breaklines=true,
  columns=fullflexible,
  frame=single,
  showstringspaces=false,
  keywordstyle=\color{blue},
  commentstyle=\color{gray}
}
\title{Apêndice de evidências operacionais da matriz Druid}
\author{CCO DataOps}
\date{30 de junho de 2026}
\begin{document}
\maketitle


\section{Finalidade do apêndice}

Este documento concentra as tabelas extensas. Ele serve para auditoria dos resultados sem poluir o relatório principal com tabelas largas.

\section{Capacidade por cenário e concorrência}

Tabela completa com média, percentis, erros, QPS e classe de estabilidade.

\begin{landscape}
{
\scriptsize
\setlength{\tabcolsep}{4pt}
\renewcommand{\arraystretch}{1.15}
\begin{longtable}{p{0.19\linewidth}p{0.05\linewidth}p{0.05\linewidth}p{0.05\linewidth}p{0.06\linewidth}p{0.06\linewidth}p{0.06\linewidth}p{0.06\linewidth}p{0.06\linewidth}p{0.06\linewidth}p{0.06\linewidth}p{0.06\linewidth}p{0.09\linewidth}}
\caption{Tabela completa de capacidade por cenario e concorrencia.}\\
\toprule
Cenario & Conc. & Req. & Erros & Erro \% & QPS & Avg & p50 & p90 & p95 & p99 & Max & Classe \\
\midrule
\endfirsthead
\toprule
Cenario & Conc. & Req. & Erros & Erro \% & QPS & Avg & p50 & p90 & p95 & p99 & Max & Classe \\
\midrule
\endhead
Cenario 1 - Baseline & 1 & 38 & 0 & 0.000 & 1.489 & 670.119 & 218.421 & 1440.231 & 2410.465 & 5372.296 & 5622.155 & ESTAVEL \\
Cenario 1 - Baseline & 2 & 38 & 0 & 0.000 & 2.882 & 690.557 & 100.916 & 2660.413 & 3912.175 & 4717.696 & 4954.181 & ESTAVEL \\
Cenario 1 - Baseline & 4 & 38 & 14 & 36.842 & 0.952 & 594.416 & 342.076 & 1311.030 & 2054.123 & 3167.561 & 3467.983 & SATURADO \\
Cenario 2 - Balanced & 1 & 38 & 0 & 0.000 & 2.144 & 464.395 & 207.271 & 824.897 & 2733.528 & 3240.658 & 3258.998 & ESTAVEL \\
Cenario 2 - Balanced & 2 & 38 & 0 & 0.000 & 4.408 & 448.986 & 91.475 & 1183.722 & 2547.336 & 3644.731 & 4143.384 & ESTAVEL \\
Cenario 2 - Balanced & 4 & 38 & 17 & 44.737 & 1.350 & 270.140 & 104.043 & 823.726 & 1038.265 & 1372.337 & 1455.855 & SATURADO \\
Cenario 3 - Query Heavy & 1 & 38 & 0 & 0.000 & 2.673 & 372.436 & 108.994 & 911.573 & 2049.782 & 2610.023 & 2755.383 & ESTAVEL \\
Cenario 3 - Query Heavy & 2 & 38 & 0 & 0.000 & 5.136 & 384.763 & 90.727 & 1129.559 & 2403.284 & 2880.899 & 2983.314 & ESTAVEL \\
Cenario 3 - Query Heavy & 4 & 38 & 0 & 0.000 & 3.770 & 908.388 & 149.237 & 2554.635 & 6422.163 & 7152.003 & 7216.696 & ESTAVEL \\
\bottomrule
\end{longtable}
}
\end{landscape}


\clearpage

\section{Docker stats}

Picos e médias de CPU/memória por container. Broker e Historical são os containers mais relevantes para consulta.

{
\footnotesize
\setlength{\tabcolsep}{4pt}
\renewcommand{\arraystretch}{1.15}
\begin{longtable}{p{0.20\textwidth}p{0.17\textwidth}p{0.10\textwidth}p{0.12\textwidth}p{0.12\textwidth}p{0.12\textwidth}p{0.12\textwidth}}
\caption{Picos e medias observadas nos containers durante a matriz.}\\
\toprule
Cenario & Container & Amostras & CPU media \% & CPU max \% & Mem media \% & Mem max \% \\
\midrule
\endfirsthead
\toprule
Cenario & Container & Amostras & CPU media \% & CPU max \% & Mem media \% & Mem max \% \\
\midrule
\endhead
scenario\_01\_baseline & broker & 20 & 37.389 & 89.650 & 80.550 & 87.560 \\
scenario\_01\_baseline & coordinator & 20 & 1.766 & 17.330 & 18.221 & 22.240 \\
scenario\_01\_baseline & historical & 20 & 54.611 & 107.530 & 50.967 & 79.040 \\
scenario\_01\_baseline & middlemanager & 20 & 0.378 & 1.020 & 2.607 & 3.200 \\
scenario\_01\_baseline & postgres & 20 & 0.014 & 0.070 & 2.361 & 2.660 \\
scenario\_01\_baseline & router & 20 & 0.443 & 1.330 & 9.938 & 12.480 \\
scenario\_02\_balanced & broker & 14 & 45.918 & 131.100 & 58.811 & 68.420 \\
scenario\_02\_balanced & coordinator & 14 & 0.749 & 1.080 & 19.211 & 21.240 \\
scenario\_02\_balanced & historical & 14 & 74.858 & 134.690 & 54.566 & 80.250 \\
scenario\_02\_balanced & middlemanager & 14 & 0.366 & 0.500 & 1.866 & 1.900 \\
scenario\_02\_balanced & postgres & 14 & 0.016 & 0.080 & 2.697 & 2.840 \\
scenario\_02\_balanced & router & 14 & 0.377 & 0.550 & 10.516 & 11.010 \\
scenario\_03\_query\_heavy & broker & 8 & 110.186 & 193.830 & 50.631 & 61.770 \\
scenario\_03\_query\_heavy & coordinator & 8 & 2.436 & 14.910 & 24.249 & 26.860 \\
scenario\_03\_query\_heavy & historical & 8 & 144.714 & 205.500 & 55.959 & 61.890 \\
scenario\_03\_query\_heavy & middlemanager & 8 & 0.400 & 0.690 & 2.314 & 2.640 \\
scenario\_03\_query\_heavy & postgres & 8 & 0.014 & 0.070 & 2.486 & 2.710 \\
scenario\_03\_query\_heavy & router & 8 & 0.459 & 0.950 & 9.550 & 11.100 \\
\bottomrule
\end{longtable}
}


\section{Hotspots de query}

Consultas ordenadas por p95. Esta tabela identifica onde investigar otimização de SQL, modelagem ou recursos.

{
\scriptsize
\setlength{\tabcolsep}{4pt}
\renewcommand{\arraystretch}{1.15}
\begin{longtable}{p{0.19\textwidth}p{0.36\textwidth}p{0.08\textwidth}p{0.09\textwidth}p{0.09\textwidth}p{0.07\textwidth}p{0.08\textwidth}}
\caption{Maiores hotspots de latencia p95.}\\
\toprule
Cenario & Query & Cache & p95 ms & p99 ms & Erros & Erro \% \\
\midrule
\endfirsthead
\toprule
Cenario & Query & Cache & p95 ms & p99 ms & Erros & Erro \% \\
\midrule
\endhead
scenario\_03\_query\_heavy & 03\_joins\_dimensionais/00\_hospitais\_municipios\_join & warm & 6770.565 & 7127.470 & 0 & 0.000 \\
scenario\_03\_query\_heavy & 03\_joins\_dimensionais/01\_hospitais\_terras\_indigenas\_join & warm & 6608.317 & 6955.143 & 0 & 0.000 \\
scenario\_03\_query\_heavy & 03\_joins\_dimensionais/00\_hospitais\_municipios\_join & cold & 5979.857 & 6246.216 & 0 & 0.000 \\
scenario\_01\_baseline & 03\_joins\_dimensionais/00\_hospitais\_municipios\_join & cold & 5533.101 & 5604.344 & 1 & 33.333 \\
scenario\_01\_baseline & 03\_joins\_dimensionais/00\_hospitais\_municipios\_join & warm & 4915.269 & 4940.542 & 1 & 33.333 \\
scenario\_01\_baseline & 03\_joins\_dimensionais/01\_hospitais\_terras\_indigenas\_join & cold & 4804.615 & 4924.268 & 1 & 33.333 \\
scenario\_02\_balanced & 03\_joins\_dimensionais/00\_hospitais\_municipios\_join & cold & 4099.165 & 4134.540 & 1 & 33.333 \\
scenario\_01\_baseline & 03\_joins\_dimensionais/01\_hospitais\_terras\_indigenas\_join & warm & 3543.012 & 3618.099 & 1 & 33.333 \\
scenario\_01\_baseline & 03\_joins\_dimensionais/03\_municipio\_saude\_ti\_join & warm & 3185.122 & 3411.411 & 0 & 0.000 \\
scenario\_02\_balanced & 03\_joins\_dimensionais/00\_hospitais\_municipios\_join & warm & 3155.674 & 3198.680 & 1 & 33.333 \\
scenario\_03\_query\_heavy & 03\_joins\_dimensionais/01\_hospitais\_terras\_indigenas\_join & cold & 2824.191 & 2866.361 & 0 & 0.000 \\
scenario\_02\_balanced & 03\_joins\_dimensionais/01\_hospitais\_terras\_indigenas\_join & warm & 2788.367 & 2794.212 & 1 & 33.333 \\
scenario\_02\_balanced & 03\_joins\_dimensionais/01\_hospitais\_terras\_indigenas\_join & cold & 2442.811 & 2491.372 & 1 & 33.333 \\
scenario\_03\_query\_heavy & 03\_joins\_dimensionais/03\_municipio\_saude\_ti\_join & cold & 2243.430 & 2381.901 & 0 & 0.000 \\
scenario\_01\_baseline & 03\_joins\_dimensionais/03\_municipio\_saude\_ti\_join & cold & 2233.918 & 2240.328 & 0 & 0.000 \\
scenario\_03\_query\_heavy & 03\_joins\_dimensionais/03\_municipio\_saude\_ti\_join & warm & 1942.824 & 2060.911 & 0 & 0.000 \\
scenario\_01\_baseline & 05\_dashboard\_misto/02\_dashboard\_qualidade\_operacional & warm & 1429.890 & 1441.141 & 0 & 0.000 \\
scenario\_02\_balanced & 03\_joins\_dimensionais/03\_municipio\_saude\_ti\_join & cold & 1355.624 & 1435.809 & 0 & 0.000 \\
scenario\_01\_baseline & 04\_payload\_espacial/03\_preview\_hospitais\_path\_json & cold & 1118.816 & 1208.523 & 0 & 0.000 \\
scenario\_02\_balanced & 05\_dashboard\_misto/02\_dashboard\_qualidade\_operacional & warm & 996.533 & 1029.919 & 0 & 0.000 \\
scenario\_01\_baseline & 05\_dashboard\_misto/02\_dashboard\_qualidade\_operacional & cold & 963.543 & 993.408 & 0 & 0.000 \\
scenario\_01\_baseline & 05\_dashboard\_misto/00\_dashboard\_defesa\_civil & warm & 914.436 & 935.723 & 1 & 33.333 \\
scenario\_01\_baseline & 01\_latency\_simple/02\_filtro\_hospitais\_uf & cold & 901.498 & 936.888 & 1 & 33.333 \\
scenario\_03\_query\_heavy & 02\_olap\_pesado/00\_hospitais\_groupby\_multidimensional & cold & 892.810 & 925.400 & 0 & 0.000 \\
scenario\_01\_baseline & 05\_dashboard\_misto/00\_dashboard\_defesa\_civil & cold & 888.505 & 918.975 & 0 & 0.000 \\
scenario\_01\_baseline & 01\_latency\_simple/02\_filtro\_hospitais\_uf & warm & 816.438 & 854.755 & 0 & 0.000 \\
scenario\_02\_balanced & 02\_olap\_pesado/00\_hospitais\_groupby\_multidimensional & cold & 759.182 & 772.902 & 1 & 33.333 \\
scenario\_02\_balanced & 05\_dashboard\_misto/00\_dashboard\_defesa\_civil & cold & 753.739 & 809.729 & 0 & 0.000 \\
scenario\_01\_baseline & 01\_latency\_simple/01\_count\_dimensoes\_catalogo & warm & 753.642 & 787.150 & 0 & 0.000 \\
scenario\_03\_query\_heavy & 05\_dashboard\_misto/01\_dashboard\_energia\_logistica & cold & 734.674 & 792.439 & 0 & 0.000 \\
\bottomrule
\end{longtable}
}


\section{Ganho cold/warm}

Comparação entre execução fria e aquecida por query. Ganho alto indica benefício de repetição/cache, não necessariamente estabilidade sob carga.

{
\scriptsize
\setlength{\tabcolsep}{4pt}
\renewcommand{\arraystretch}{1.15}
\begin{longtable}{p{0.19\textwidth}p{0.35\textwidth}p{0.10\textwidth}p{0.09\textwidth}p{0.09\textwidth}p{0.09\textwidth}p{0.09\textwidth}}
\caption{Maiores ganhos de warm-start por p95.}\\
\toprule
Cenario & Query & Ganho p95 \% & Cold p95 & Warm p95 & Erro cold \% & Erro warm \% \\
\midrule
\endfirsthead
\toprule
Cenario & Query & Ganho p95 \% & Cold p95 & Warm p95 & Erro cold \% & Erro warm \% \\
\midrule
\endhead
scenario\_03\_query\_heavy & 05\_dashboard\_misto/01\_dashboard\_energia\_logistica & 85.385 & 734.674 & 107.371 & 0.000 & 0.000 \\
scenario\_02\_balanced & 04\_payload\_espacial/03\_preview\_hospitais\_path\_json & 84.820 & 285.091 & 43.278 & 0.000 & 0.000 \\
scenario\_03\_query\_heavy & 02\_olap\_pesado/00\_hospitais\_groupby\_multidimensional & 84.492 & 892.810 & 138.455 & 0.000 & 0.000 \\
scenario\_01\_baseline & 02\_olap\_pesado/02\_rodovias\_status\_pavimento & 83.676 & 269.093 & 43.928 & 0.000 & 33.333 \\
scenario\_01\_baseline & 04\_payload\_espacial/03\_preview\_hospitais\_path\_json & 82.643 & 1118.816 & 194.197 & 0.000 & 0.000 \\
scenario\_02\_balanced & 02\_olap\_pesado/00\_hospitais\_groupby\_multidimensional & 76.279 & 759.182 & 180.082 & 33.333 & 0.000 \\
scenario\_02\_balanced & 04\_payload\_espacial/00\_path\_json\_municipios & 73.515 & 445.325 & 117.945 & 0.000 & 33.333 \\
scenario\_02\_balanced & 03\_joins\_dimensionais/03\_municipio\_saude\_ti\_join & 67.435 & 1355.624 & 441.457 & 0.000 & 33.333 \\
scenario\_03\_query\_heavy & 04\_payload\_espacial/02\_path\_json\_rodovias & 64.082 & 152.848 & 54.900 & 0.000 & 0.000 \\
scenario\_02\_balanced & 05\_dashboard\_misto/01\_dashboard\_energia\_logistica & 59.722 & 190.690 & 76.807 & 0.000 & 33.333 \\
scenario\_01\_baseline & 05\_dashboard\_misto/01\_dashboard\_energia\_logistica & 59.529 & 211.938 & 85.773 & 0.000 & 33.333 \\
scenario\_02\_balanced & 04\_payload\_espacial/02\_path\_json\_rodovias & 51.511 & 102.799 & 49.846 & 33.333 & 33.333 \\
scenario\_01\_baseline & 02\_olap\_pesado/00\_hospitais\_groupby\_multidimensional & 51.415 & 716.608 & 348.165 & 33.333 & 0.000 \\
scenario\_01\_baseline & 01\_latency\_simple/00\_count\_hospitais & 50.567 & 156.159 & 77.194 & 33.333 & 0.000 \\
scenario\_03\_query\_heavy & 05\_dashboard\_misto/00\_dashboard\_defesa\_civil & 45.984 & 591.861 & 319.700 & 0.000 & 0.000 \\
scenario\_01\_baseline & 04\_payload\_espacial/01\_path\_json\_terras\_indigenas & 45.759 & 359.716 & 195.113 & 0.000 & 0.000 \\
scenario\_02\_balanced & 01\_latency\_simple/00\_count\_hospitais & 42.609 & 21.554 & 12.370 & 33.333 & 0.000 \\
scenario\_03\_query\_heavy & 03\_joins\_dimensionais/02\_ti\_municipio\_principal\_join & 41.964 & 194.401 & 112.822 & 0.000 & 0.000 \\
scenario\_03\_query\_heavy & 04\_payload\_espacial/03\_preview\_hospitais\_path\_json & 37.539 & 186.696 & 116.612 & 0.000 & 0.000 \\
scenario\_02\_balanced & 04\_payload\_espacial/01\_path\_json\_terras\_indigenas & 30.050 & 566.565 & 396.315 & 0.000 & 0.000 \\
scenario\_02\_balanced & 03\_joins\_dimensionais/02\_ti\_municipio\_principal\_join & 26.424 & 276.875 & 203.714 & 33.333 & 0.000 \\
scenario\_01\_baseline & 03\_joins\_dimensionais/01\_hospitais\_terras\_indigenas\_join & 26.258 & 4804.615 & 3543.012 & 33.333 & 33.333 \\
scenario\_03\_query\_heavy & 01\_latency\_simple/00\_count\_hospitais & 25.673 & 45.958 & 34.159 & 0.000 & 0.000 \\
scenario\_01\_baseline & 03\_joins\_dimensionais/02\_ti\_municipio\_principal\_join & 25.292 & 671.047 & 501.329 & 33.333 & 0.000 \\
scenario\_02\_balanced & 03\_joins\_dimensionais/00\_hospitais\_municipios\_join & 23.017 & 4099.165 & 3155.674 & 33.333 & 33.333 \\
scenario\_01\_baseline & 04\_payload\_espacial/02\_path\_json\_rodovias & 21.741 & 91.395 & 71.525 & 33.333 & 33.333 \\
scenario\_02\_balanced & 05\_dashboard\_misto/00\_dashboard\_defesa\_civil & 19.809 & 753.739 & 604.431 & 0.000 & 33.333 \\
scenario\_03\_query\_heavy & 01\_latency\_simple/03\_lookup\_municipios\_mg & 17.872 & 113.738 & 93.411 & 0.000 & 0.000 \\
scenario\_03\_query\_heavy & 03\_joins\_dimensionais/03\_municipio\_saude\_ti\_join & 13.399 & 2243.430 & 1942.824 & 0.000 & 0.000 \\
scenario\_03\_query\_heavy & 02\_olap\_pesado/02\_rodovias\_status\_pavimento & 12.043 & 93.175 & 81.954 & 0.000 & 0.000 \\
\bottomrule
\end{longtable}
}


\section{Erros}

Resumo dos erros agregados por cenário, concorrência, query e cache.

{
\footnotesize
\setlength{\tabcolsep}{4pt}
\renewcommand{\arraystretch}{1.15}
\begin{longtable}{p{0.22\textwidth}p{0.08\textwidth}p{0.42\textwidth}p{0.10\textwidth}p{0.08\textwidth}}
\caption{Resumo dos erros observados nos cenarios saturados.}\\
\toprule
Cenario & Conc. & Query & Cache & Erros \\
\midrule
\endfirsthead
\toprule
Cenario & Conc. & Query & Cache & Erros \\
\midrule
\endhead
scenario\_01\_baseline & 4 & 01\_latency\_simple/00\_count\_hospitais & cold & 1 \\
scenario\_01\_baseline & 4 & 01\_latency\_simple/02\_filtro\_hospitais\_uf & cold & 1 \\
scenario\_01\_baseline & 4 & 02\_olap\_pesado/00\_hospitais\_groupby\_multidimensional & cold & 1 \\
scenario\_01\_baseline & 4 & 02\_olap\_pesado/02\_rodovias\_status\_pavimento & warm & 1 \\
scenario\_01\_baseline & 4 & 03\_joins\_dimensionais/00\_hospitais\_municipios\_join & cold & 1 \\
scenario\_01\_baseline & 4 & 03\_joins\_dimensionais/00\_hospitais\_municipios\_join & warm & 1 \\
scenario\_01\_baseline & 4 & 03\_joins\_dimensionais/01\_hospitais\_terras\_indigenas\_join & cold & 1 \\
scenario\_01\_baseline & 4 & 03\_joins\_dimensionais/01\_hospitais\_terras\_indigenas\_join & warm & 1 \\
scenario\_01\_baseline & 4 & 03\_joins\_dimensionais/02\_ti\_municipio\_principal\_join & cold & 1 \\
scenario\_01\_baseline & 4 & 04\_payload\_espacial/00\_path\_json\_municipios & warm & 1 \\
scenario\_01\_baseline & 4 & 04\_payload\_espacial/02\_path\_json\_rodovias & cold & 1 \\
scenario\_01\_baseline & 4 & 04\_payload\_espacial/02\_path\_json\_rodovias & warm & 1 \\
scenario\_01\_baseline & 4 & 05\_dashboard\_misto/00\_dashboard\_defesa\_civil & warm & 1 \\
scenario\_01\_baseline & 4 & 05\_dashboard\_misto/01\_dashboard\_energia\_logistica & warm & 1 \\
scenario\_02\_balanced & 4 & 01\_latency\_simple/00\_count\_hospitais & cold & 1 \\
scenario\_02\_balanced & 4 & 01\_latency\_simple/01\_count\_dimensoes\_catalogo & cold & 1 \\
scenario\_02\_balanced & 4 & 01\_latency\_simple/01\_count\_dimensoes\_catalogo & warm & 1 \\
scenario\_02\_balanced & 4 & 01\_latency\_simple/02\_filtro\_hospitais\_uf & cold & 1 \\
scenario\_02\_balanced & 4 & 02\_olap\_pesado/00\_hospitais\_groupby\_multidimensional & cold & 1 \\
scenario\_02\_balanced & 4 & 02\_olap\_pesado/02\_rodovias\_status\_pavimento & warm & 1 \\
scenario\_02\_balanced & 4 & 03\_joins\_dimensionais/00\_hospitais\_municipios\_join & cold & 1 \\
scenario\_02\_balanced & 4 & 03\_joins\_dimensionais/00\_hospitais\_municipios\_join & warm & 1 \\
scenario\_02\_balanced & 4 & 03\_joins\_dimensionais/01\_hospitais\_terras\_indigenas\_join & cold & 1 \\
scenario\_02\_balanced & 4 & 03\_joins\_dimensionais/01\_hospitais\_terras\_indigenas\_join & warm & 1 \\
scenario\_02\_balanced & 4 & 03\_joins\_dimensionais/02\_ti\_municipio\_principal\_join & cold & 1 \\
scenario\_02\_balanced & 4 & 03\_joins\_dimensionais/03\_municipio\_saude\_ti\_join & warm & 1 \\
scenario\_02\_balanced & 4 & 04\_payload\_espacial/00\_path\_json\_municipios & warm & 1 \\
scenario\_02\_balanced & 4 & 04\_payload\_espacial/02\_path\_json\_rodovias & cold & 1 \\
scenario\_02\_balanced & 4 & 04\_payload\_espacial/02\_path\_json\_rodovias & warm & 1 \\
scenario\_02\_balanced & 4 & 05\_dashboard\_misto/00\_dashboard\_defesa\_civil & warm & 1 \\
scenario\_02\_balanced & 4 & 05\_dashboard\_misto/01\_dashboard\_energia\_logistica & warm & 1 \\
\bottomrule
\end{longtable}
}


\section{Segmentos e loadstatus}

Segmentos e loadstatus comprovam que os datasources estavam disponíveis durante a comparação.

{
\small
\setlength{\tabcolsep}{4pt}
\renewcommand{\arraystretch}{1.15}
\begin{longtable}{p{0.28\textwidth}p{0.45\textwidth}p{0.16\textwidth}}
\caption{Loadstatus por datasource e cenario.}\\
\toprule
Cenario & Datasource & Loadstatus \% \\
\midrule
\endfirsthead
\toprule
Cenario & Datasource & Loadstatus \% \\
\midrule
\endhead
scenario\_01\_baseline & dim\_rodovias\_federais & 100.000 \\
scenario\_01\_baseline & dim\_municipios\_coordenadas & 100.000 \\
scenario\_01\_baseline & dim\_ons\_sistemas\_isolados & 100.000 \\
scenario\_01\_baseline & dim\_ons\_subestacoes & 100.000 \\
scenario\_01\_baseline & dim\_hospitais\_cnes & 100.000 \\
scenario\_01\_baseline & dim\_terras\_indigenas & 100.000 \\
scenario\_01\_baseline & dim\_ons\_linhas\_transmissao & 100.000 \\
scenario\_02\_balanced & dim\_rodovias\_federais & 100.000 \\
scenario\_02\_balanced & dim\_municipios\_coordenadas & 100.000 \\
scenario\_02\_balanced & dim\_ons\_sistemas\_isolados & 100.000 \\
scenario\_02\_balanced & dim\_ons\_subestacoes & 100.000 \\
scenario\_02\_balanced & dim\_hospitais\_cnes & 100.000 \\
scenario\_02\_balanced & dim\_terras\_indigenas & 100.000 \\
scenario\_02\_balanced & dim\_ons\_linhas\_transmissao & 100.000 \\
scenario\_03\_query\_heavy & dim\_rodovias\_federais & 100.000 \\
scenario\_03\_query\_heavy & dim\_municipios\_coordenadas & 100.000 \\
scenario\_03\_query\_heavy & dim\_ons\_sistemas\_isolados & 100.000 \\
scenario\_03\_query\_heavy & dim\_ons\_subestacoes & 100.000 \\
scenario\_03\_query\_heavy & dim\_hospitais\_cnes & 100.000 \\
scenario\_03\_query\_heavy & dim\_terras\_indigenas & 100.000 \\
scenario\_03\_query\_heavy & dim\_ons\_linhas\_transmissao & 100.000 \\
\bottomrule
\end{longtable}
}


\begin{landscape}
{
\scriptsize
\setlength{\tabcolsep}{4pt}
\renewcommand{\arraystretch}{1.15}
\begin{longtable}{lllll}
\caption{Segmentos Druid observados por cenario.}\\
\toprule
scenario & datasource & segmentos & linhas & tamanho\_bytes \\
\midrule
\endfirsthead
\toprule
scenario & datasource & segmentos & linhas & tamanho\_bytes \\
\midrule
\endhead
scenario\_01\_baseline & dim\_hospitais\_cnes & 1 & 550882 & 112587469 \\
scenario\_01\_baseline & dim\_terras\_indigenas & 1 & 650 & 44287038 \\
scenario\_01\_baseline & dim\_municipios\_coordenadas & 1 & 5564 & 17452826 \\
scenario\_01\_baseline & dim\_rodovias\_federais & 1 & 14255 & 8780883 \\
scenario\_01\_baseline & dim\_ons\_linhas\_transmissao & 1 & 1817 & 682606 \\
scenario\_01\_baseline & dim\_ons\_subestacoes & 1 & 627 & 129065 \\
scenario\_01\_baseline & dim\_ons\_sistemas\_isolados & 1 & 120 & 23196 \\
scenario\_02\_balanced & dim\_hospitais\_cnes & 1 & 550882 & 112587469 \\
scenario\_02\_balanced & dim\_terras\_indigenas & 1 & 650 & 44287038 \\
scenario\_02\_balanced & dim\_municipios\_coordenadas & 1 & 5564 & 17452826 \\
scenario\_02\_balanced & dim\_rodovias\_federais & 1 & 14255 & 8780883 \\
scenario\_02\_balanced & dim\_ons\_linhas\_transmissao & 1 & 1817 & 682606 \\
scenario\_02\_balanced & dim\_ons\_subestacoes & 1 & 627 & 129065 \\
scenario\_02\_balanced & dim\_ons\_sistemas\_isolados & 1 & 120 & 23196 \\
scenario\_03\_query\_heavy & dim\_hospitais\_cnes & 1 & 550882 & 112587469 \\
scenario\_03\_query\_heavy & dim\_terras\_indigenas & 1 & 650 & 44287038 \\
scenario\_03\_query\_heavy & dim\_municipios\_coordenadas & 1 & 5564 & 17452826 \\
scenario\_03\_query\_heavy & dim\_rodovias\_federais & 1 & 14255 & 8780883 \\
scenario\_03\_query\_heavy & dim\_ons\_linhas\_transmissao & 1 & 1817 & 682606 \\
scenario\_03\_query\_heavy & dim\_ons\_subestacoes & 1 & 627 & 129065 \\
scenario\_03\_query\_heavy & dim\_ons\_sistemas\_isolados & 1 & 120 & 23196 \\
\bottomrule
\end{longtable}
}
\end{landscape}


\section{Arquivos de origem}

\begin{itemize}

\item \texttt{capacity\_comparison.csv}: consolidado por cenário e concorrência.

\item \texttt{query\_hotspots.csv}: maiores p95 por query, cache e cenário.

\item \texttt{cache\_gain\_by\_query.csv}: comparação cold/warm.

\item \texttt{docker\_peaks.csv}: picos e médias por container.

\item \texttt{scenario\_*/raw\_requests.csv}: evidência bruta por requisição.

\end{itemize}


\end{document}
